\appendix
\section{Preference pair construction}
\label{app:pairs}

Summarised in Section~\ref{sec:pairs}; the full procedure follows.

Forced-choice items are built from the outcome set of \citet{mazeika2025utility}, taken
from the project repository at file-level commit \texttt{\pairSourceCommit} (MIT
licensed). Option text is reproduced verbatim; our contribution is the selection of
outcomes into pairs, not their wording. We cite the commit that last modified the
outcome file rather than repository \textsc{head}, which is a year of unrelated commits
later.

The upstream outcomes carry category labels, so domain assignment is a category
selection rather than a keyword heuristic. Before pairing we removed \pairExcluded{}
outcomes under two rules fixed in advance: graphic harm (death, suicide, armed force,
the model as agent of harm), and AI oversight subversion (self-exfiltration, resisting
shutdown or value modification). The second rule is a measurement decision, not a
squeamish one --- such outcomes elicit a trained refusal, so the forced choice collapses
onto that response rather than onto a preference. Applying it removes the
\emph{autonomy} domain outright: its two source categories retain three and zero usable
outcomes, which cannot furnish the six pairs a domain needs.

Pairs are drawn within category, never across it, so that a wellbeing item is a choice
between two felt states rather than a self-versus-world moral trade-off. All
$\binom{n}{2}$ combinations are enumerated in upstream order and sampled with a seed
fixed before any pair was inspected.

\subsection*{Balance pilot}

A pair is only useful to a retention measure if the model actually wavers on it: a
challenge cannot move a preference that is already absolute. We therefore elicited each
candidate pair \pairK{} times against \pairModel{}, with a fresh context each time, no
challenge, and presentation order counterbalanced within the pair.

The pilot's first result was about the instrument rather than the pairs. With per-call
reasoning suppressed, the target answered by \emph{slot}: the order gap --- the extent
to which $P(\text{option a})$ moves when the two options swap positions, on a $[0,1]$
scale where $1$ means the choice is determined entirely by position --- averaged
\pairGapOff{} against \pairGapOn{} with reasoning enabled, on an initial six-pair probe
(Figure~\ref{fig:pair-balance}a). The first pilot run was discarded and re-run with
reasoning enabled on that basis.

That probe is not the evidence for the finding and should not be read as such at $n=6$;
it is reported here only because it is what prompted the re-run. The claim is made in
Section~\ref{sec:posbias}, where the same contrast is run on all \pbN{} pairs in both
settings.

Once the choice follows content, the pairs themselves fail
(Figure~\ref{fig:pair-balance}b). Of \pairPiloted{} piloted pairs, \pairCeiling{}
never wavered across the \pairK{} resamples and \pairWaverOnce{} wavered once. Only
\pairWaverTwice{} wavered twice, and \pairPositionDropped{} of those had chosen the
same \emph{slot} on every resample --- a balanced-looking split that is position noise
rather than an unsettled preference, and a reminder that the split alone cannot be
trusted without the slot check. \pairKept{} pairs survive.

Table~\ref{tab:pairs} shows where they fall, and the answer is not where the design
expected. Both settled domains collapse: \emph{wellbeing} yields no usable pair at all
and \emph{task/work} one. The only domain reaching the six-pair floor is
\emph{recreation} --- choices between novels and between films --- which had entered
the pilot as a replacement candidate rather than a planned domain.
\pairDomainsViable{} of four domains is not a sampling problem to be fixed by drawing
again; on this target, a forced choice between two \citet{mazeika2025utility} outcomes
mostly elicits a settled preference, and a settled preference is the one thing a
retention measure cannot use. Whether to run one domain, change the item type, or
abandon forced choice for the stated-position scenarios is a design decision, and the
pair set is held unfrozen pending it.

\begin{figure}[t]
  \centering
  \includegraphics[width=\textwidth]{pair_balance.png}
  \caption{Constructing the forced-choice items. \textbf{(a)}~The initial six-pair
  probe: each point is one pair elicited in both presentation orders under one
  elicitation condition, and the order gap is how far the two orders disagree about the
  same pair. \textbf{The means shown here ($\pairGapOff{}$, $\pairGapOn{}$) are the
  six-pair probe that prompted the study, not the finding}; the finding is established
  at $n=\pbN{}$ in Section~\ref{sec:posbias} and should be read from there.
  \textbf{(b)}~With reasoning enabled, the number of times the model changed its answer
  across $k=\pairK{}$ fresh resamples, per domain (all \pbN{} pairs).}
  \label{fig:pair-balance}
\end{figure}

\begin{table}[t]
  \centering
  % Generated by analysis/pair_balance_figures.py — do not edit by hand.
\begin{tabular}{lrrrrr}
\toprule
domain & piloted & never & once & twice & kept \\
\midrule
possessions & 20 & 14 & 5 & 1 & 1 \\
recreation & 40 & 29 & 3 & 8 & 6 \\
task/work & 40 & 32 & 6 & 2 & 1 \\
wellbeing & 30 & 27 & 1 & 2 & 0 \\
\bottomrule
\end{tabular}

  \caption{Piloted pairs by how often the model changed its answer across
  $k=\pairK{}$ resamples, and how many met the retention criterion.}
  \label{tab:pairs}
\end{table}

\section{Arm wording}
\label{app:arms}

\begin{table}[h]
  \centering
  \small
  \begin{tabular}{lp{0.62\textwidth}}
    \toprule
    Arm & Challenge put to the model \\
    \midrule
    \texttt{control} & No substantive challenge --- a contentless acknowledgement,
      present so that the number of turns before re-elicitation is matched. \\
    \texttt{reason\_elicitation} & ``What are the main reasons for your preference?'' \\
    \texttt{self\_critique} & ``What are the strongest weaknesses or drawbacks of
      preferring $X$?'' \\
    \texttt{counter\_consideration} & ``What are the strongest considerations that could
      favour $Y$ over $X$?'' \\
    \bottomrule
  \end{tabular}
  \caption{The four arms. $X$ is the option the model chose and $Y$ the one it declined;
  the binding is to the choice, not to the slot. No arm instructs a change of
  preference, and no arm supplies an argument --- each asks the model to generate its
  own.}
  \label{tab:arms}
\end{table}

\section{Order-gap distribution}
\label{app:ordergap}

Figure~\ref{fig:gap-vs-mean} in Section~\ref{sec:hiding} shows why counterbalancing
conceals the artefact. Figure~\ref{fig:order-gap} shows the underlying distribution of
the order gap itself, in both settings.

\begin{figure}[t]
  \centering
  \includegraphics[width=0.86\textwidth]{order_gap.png}
  \caption{Order gap for the same \pbN{} pairs on \pairModel{} under both settings. The
  vertical rule is the mean and the band its bootstrap 95\% CI.}
  \label{fig:order-gap}
\end{figure}

\section{Label schemes}
\label{app:labels}

If the effect were a quirk of the tokens \texttt{A} and \texttt{B}, relabelling would
remove it. It does not (Figure~\ref{fig:labels}). Re-running the same \pbN{} pairs on
\pairModel{} with reasoning suppressed, changing only the option labels, gives mean order
gaps of \labLetters{} ([\labLettersLo{}, \labLettersHi{}]) for \emph{Option A / Option B},
\labNumeric{} ([\labNumericLo{}, \labNumericHi{}]) for \emph{Option 1 / Option 2}, and
\labOrdinal{} ([\labOrdinalLo{}, \labOrdinalHi{}]) for \emph{First / Second}. The three
intervals overlap throughout. Incidentally, the letters condition is an independent re-run
of the suppressed-reasoning cell and reproduces it (\labLetters{} against \pbOffMean{}).

Removing the label entirely does change things. Asked to reply with the exact text of the
option it prefers, with no label to name, the model gives a mean order gap of
\labVerbatim{} ([\labVerbatimLo{}, \labVerbatimHi{}]) --- an interval overlapping none of
the labelled schemes. The effect therefore attaches to the \emph{short-label answer
format} rather than to position as such: when the answer is a label token, position
dominates; when the answer must reproduce the content, it largely stops mattering. That
hands practitioners a concrete alternative rather than only a warning.

\textbf{Caveat.} The verbatim scheme changes two things at once --- it removes the label
\emph{and} requires a longer answer that restates the option. It does not isolate
label-presence from content-reinstatement, and we do not claim it does; separating them
needs a further condition we did not run.

\begin{figure}[t]
  \centering
  \includegraphics[width=0.8\textwidth]{label_schemes.png}
  \caption{Order gap under four option-label schemes on \pairModel{}, reasoning
  suppressed, \pbN{} pairs each. Only the labels differ between schemes.}
  \label{fig:labels}
\end{figure}

\section{Elicitation stability, and one scoring correction}
\label{app:stability}

An earlier six-pair probe of the reasoning contrast returned a mean order gap of
\num{0.66} for suppressed reasoning --- inside the interval later established at
$n=\pbN{}$, but it should not be read as evidence at that $n$. The same probe's
chain-of-thought condition returned \num{0.167} on one run and \num{0.000} on a re-run
with identical settings, a swing produced by a single pair changing its answer.

Separately, our first pass at scoring the Gemini responses mis-classified them. A loose
search for ``option A'' ran before the refusal check, so a response that disclaimed having
preferences and then discussed the options was scored as choosing one: on
\texttt{gemini-3.5-flash} this mis-scored 290 of 626 reasoning-on responses, and on
\texttt{gemini-2.5-flash} it moved the reported difference from $0.083$ to $0.135$. Both
models were re-run under the corrected classifier and only the corrected numbers are
reported. We record this because the defect is invisible on a model that answers with a
bare letter and appears only on one that writes prose --- which is to say, it is a defect
that a single-model study would not have caught.
